% Separate post-CP25 candidate. No assertion is promoted by finite checks.
\subsection{One-slope quotients and the first tame boundary}

Throughout this subsection \(p\ge5\), \(q=p^2\), \(N=q-1\), and
\(K\) is complete discretely valued of mixed characteristic \((0,p)\),
with algebraically closed residue field, \(\mu_p\subset K\), and
absolute ramification \(e\) prime to \(p\). Write
\[
 p-1\mid e,\qquad s=e(p+1),\qquad
 L/K\text{ tame of degree }N,\qquad v_L(p)=eN=(p-1)s .
 \tag{B-RD69.1}
\]
The extension of degree \(N\) is used even when the special action
is not faithful. By Lemma~\ref{lem:B-RD49-inertia}, it is a good
field for every whole potentially supersingular height-two group.
Fix \(\sigma\pi_L=\zeta\pi_L\), with \(\zeta\) primitive of order \(N\).
An integral Honda framing and tame averaging, as in
Lemma~\ref{lem:B-RD67-normalization}, give an exact semilinear
coordinate whose special tangent multiplier is \(\zeta^c\).
The finite special character has eigenvalue exponents \(\{c,pc\}\).
No square-class normalization of its determinant is made here.

\begin{lemma}[Comparison rules conditional on unequal slopes]
\label{lem:B-RD69-rules}
\status{Candidate; complete proof supplied, independent audits pending.}
If the first-level Newton polygon has two slopes, it has a
unique high line \(H\) and an integer \(B\) with
\[
 \begin{gathered}
 0<B<pe,\qquad B\equiv c\pmod{p+1},\\
 \alpha=s-B,\qquad \beta=B/p,\qquad \alpha>\beta,\\
 \text{high-line character }\zeta^{c-B}.
 \end{gathered}
 \tag{B-RD69.2}
\]
There are \(p-1\) nonzero points of valuation \(\alpha\) and
\(q-p\) of valuation \(\beta\). Every order-\(p^2\) point whose
\(p\)-multiple is outside \(H\) has \(v_L\)-valuation \(B/p^3\).
For two such groups with isomorphic actual \(p^2\)-torsion:

\begin{enumerate}
\item If residual wild inertia is nontrivial, the high lines
agree. Two parameters prime to \(p\) are equal and give equal
special characters; exactly one parameter cannot be divisible
by \(p\).
\item If residual inertia is tame, the first-level modules are
rational over \(L\), and \(B_i=pb_i\). If both \(b_i\) are prime
to \(p\), they are equal, the high lines agree, and the special
characters agree.
\end{enumerate}
These statements impose no parity condition on \(B\) or \(c\).
\end{lemma}
\begin{proof}
The exact universal coefficient bounds and their original source
are given in Lemma~\ref{lem:B-RD49-point}:
\cite[Proposition~5.7; equations (12.3)--(12.5);
Proposition~12.10 and proof]{HopkinsGross1994Bundles}.
The only possible negative Newton vertices are
\((1,(p-1)s),(p,b),(q,0)\), where \(b=v_L(u)\).
The high line is inertia stable. Its character is
\(\mathbf F_p^\times\)-valued and hence trivial over \(L\);
its coordinates have integral valuation. Therefore
\(\alpha=s-b/(p-1)\) is integral, so \(b=(p-1)B\).
The strict slope inequality is exactly \(0<B<pe\).
Since \(s\) is a multiple of \(N\), the leading high-line
character is \(\zeta^{c+\alpha}=\zeta^{c-B}\).
It is \(\mathbf F_p\)-valued, giving the displayed congruence.
At a low target of valuation \(\beta\), the inverse Newton polygon
has one segment of degree \(q\), giving \(B/p^3\) for every
preimage.

For clarity, the common-point-field argument is rechecked without
the nonsquare hypothesis of Lemma~\ref{lem:B-RD68-rules}.
Let \(R/L\) be the common first-level division field. The high
line and the determinant are trivial over \(L\), so the residual
image there is a subgroup of the upper unipotent group of order
\(p\). Thus \(R=L\) or \([R:L]=p\).
In the latter case the unique wild fixed line is \(H\), so the
two high lines agree. A low coordinate has \(v_R\)-valuation
\(B\); its nonzero wild displacement has valuation \(p\alpha\).
If \(r\) is the actual lower break, the fixed-coefficient
uniformizer expansion of Lemma~\ref{lem:B-RD60-break} gives
\[
 r=ps-(p+1)B \quad(p\nmid B),\qquad
 r<ps-(p+1)B \quad(p\mid B).
 \tag{B-RD69.3}
\]
Explicitly, choose residue representatives in a coefficient
ring fixed by the automorphism. For a coordinate of valuation
\(B\), its first change has valuation \(B+r\) if \(p\nmid B\);
if \(p\mid B\), the first monomial changes strictly later,
as do all higher monomials. This proves the strict version too.
Two prime-to-\(p\) parameters give the same \(B\). Their common
high-line character then identifies their \(c\)'s modulo \(N\).

If \(p\nmid B_1\) and \(p\mid B_2\), (B-RD69.3) gives
\(B_1>B_2\). Choose corresponding primitive \(P_i\) whose
\(p\)-multiples are outside the common high line. Their
\(v_R\)-valuations are \(B_i/p^2\). Their stabilizers under
the actual torsion isomorphism agree, so their point fields
\(E/R\) agree. Such a point field is a torsor under the
\(R\)-rational first-level module, of degree at most \(p^2\).
The first valuation denominator forces degree \(p^2\).
Every first-level translation therefore occurs, and the
normalized valuations in \(E\) are \(B_i\).
A translation outside \(H\) has displacement valuation \(p^2B_i\).
The first coordinate gives its exact break \(NB_1\), while
the universal lower bound for the other gives
\[
 p^2B_2\ge B_2+NB_1,
 \tag{B-RD69.4}
\]
contradicting \(B_1>B_2\).

If residual inertia is tame, the residual division field over
\(L\) is trivial, hence \(B_i=pb_i\). Choose a common residual
vector outside the two high lines; such a vector exists because
two lines do not cover \(\mathbf F_p^2\).
When both \(b_i\) are prime to \(p\), the corresponding primitive
points have valuations \(b_i/p^2\), and their common point field
over \(L\) has degree \(p^2\). A translation outside both high
lines fixes the actual break as \(Nb_1=Nb_2\).
Thus \(b_1=b_2=b\). If the high lines differed, a translation
high in one model and low in the other would have different
displacement valuations. Both coordinates have the same
prime-to-\(p\) valuation \(b\), however, so each changes in
valuation \(b+r\) for that automorphism. This contradiction
identifies the high lines; their characters identify \(c\).
\end{proof}

\begin{lemma}[Actual quotients, including the equal-slope case]
\label{lem:B-RD69-quotients}
\status{Candidate; complete proof supplied, independent audits pending.}
The quotient of a two-slope group by its high line has parameter
\[
 B_H=pB\quad(B<e),\qquad B_H=s-B\quad(B>e).
 \tag{B-RD69.5}
\]
At \(B=e\), its quotient instead has one slope: all its nonzero
\(p\)-torsion points have valuation \(e\).
If the first-level module is \(L\)-rational and \(J\) is an
inertia-stable low line, its quotient has parameter \(B_J=B/p\).
If the original group has one slope, all its nonzero
\(p\)-torsion points have valuation \(e\), and its quotient by
any inertia-stable rank-\(p\) line has two slopes
\[
 \alpha'=pe,\qquad\beta'=e/p,\qquad B'=e.
 \tag{B-RD69.6}
\]
All these are actual good formal quotients over \(\mathcal O_L\).
Their generic groups descend to \(K\), their special tangent
exponents are \(pc\), and their finite special characters are
unchanged. An actual \(p^3\)-torsion isomorphism preserving the
chosen lines induces an actual \(p^2\)-torsion isomorphism of
the quotients.
\end{lemma}
\begin{proof}
Take the finite flat closure of the chosen line in the good
model. Its quotient is a good one-dimensional formal group;
the quotient coordinate map has the form
\[
 T\longmapsto U(T)\prod_{Q\in J}(T-T(Q)),\qquad U(T)\in
 \mathcal O_L[[T]]^\times .
 \tag{B-RD69.7}
\]
This is Weierstrass preparation for the actual finite-flat
isogeny. In particular all asserted valuations are sums of
actual difference valuations.

For the unequal polygon, quotienting by the high line multiplies
every valuation below \(\alpha\) by \(p\).
The image of \(M[p]/H\) has valuation \(B\).
If \(B>e\), the inverse polygon over a nonzero high point has
\(q\) roots of valuation \(\alpha/q\); their quotient images
give \(q-p\) points of valuation \(\alpha/p\), while the
remaining \(p-1\) points have valuation \(B>\alpha/p\).
This gives \(B_H=s-B\).
If \(B<e\), the inverse roots over each nonzero high target
have valuations \((s-pB)/p\), \(p\) times, and \(B/p\),
\(q-p\) times. All lie below \(\alpha\). Including the image
of \(M[p]/H\) and dividing preimage counts by \(p\) gives
\(p-1\) points of valuation \(s-pB\) and \(q-p\) of valuation
\(B\), hence \(B_H=pB\).
At the equality \(B=e\), the inverse polygon has one segment
of slope \(-e/p\), with \(q\) roots. Both contributions to
the quotient therefore have valuation \(e\).

For a low line every nonzero class of \(M[p]/J\) has a unique
high representative. Its quotient valuation is
\(\alpha+(p-1)\beta=s-B/p\).
Preimages of nonzero low targets have valuation \(\beta/q\);
after quotienting, the \(q-p\) resulting points have valuation
\(\beta/p=B/p^2\). This proves \(B_J=B/p\), an integer because
the first-level module is \(L\)-rational.

Now start from a one-slope polygon. The product of the \(q-1\)
nonzero \(p\)-torsion parameters has valuation \(v_L(p)=eN\),
so every such parameter has valuation \(e\). For a nonzero
class of \(M[p]/J\), all \(p\) differences in (B-RD69.7) are
nonzero \(p\)-torsion parameters and have valuation \(e\).
Its image therefore has valuation \(pe\), giving \(p-1\)
nonzero quotient points. A nonzero target in \(J\) has valuation
\(e\); the inverse Newton polygon is one segment, with \(q\)
roots of valuation \(e/q\). These values are strictly below
those of nonzero points of \(J\). Formula (B-RD69.7) then gives
valuation \(e/p\). There are \((p-1)q/p=q-p\) distinct such
quotient points. The two contributions are disjoint and exhaust
the \(q-1\) nonzero points, proving (B-RD69.6).

The unique rank-\(p\) subgroup of the one-dimensional connected
special formal group is its Frobenius kernel. Thus the special
quotient is its Frobenius twist; the inherited special tangent
exponent is \(pc\), and the unordered pair \(\{c,pc\}\) is
unchanged. Generic inertia stability of \(J\) gives actual
generic descent of the quotient to \(K\).
More explicitly, Tate full faithfulness extends the original
semilinear generic descent maps and their inverses to the good
models. They preserve the schematic closure of \(J\), hence
induce semilinear quotient maps satisfying the same cocycle.
Only the generic quotient is being descended through the
ramified extension; no integral good model over \(\mathcal O_K\)
is asserted.
Finally the actual generic identity
\[
 (M/J)[p^2]=\{P\in M[p^3]:p^2P\in J\}/J
 \tag{B-RD69.8}
\]
proves the claimed passage from a \(p^3\)-comparison to a
\(p^2\)-comparison. Here both sides are finite groups over
the characteristic-zero field, so the identity also identifies
their finite group schemes and their Galois actions.
\end{proof}

\begin{proposition}[Excluding a mixture of Newton polygons]
\label{prop:B-RD69-one-slope}
\status{Candidate; complete proof supplied, independent audits pending.}
For every \(e\) as in (B-RD69.1), two whole potentially
supersingular height-two groups with isomorphic actual
\(p^3\)-torsion cannot have one one-slope and one two-slope
first-level Newton polygon. If both have one slope, their finite
special characters agree.
\end{proposition}
\begin{proof}
For a one-slope group the map taking a \(p\)-torsion coordinate
to the leading coefficient of \(T/\pi_L^e\) is injective and
\(\mathbf F_p\)-linear. Indeed every nonzero difference is a
nonzero \(p\)-torsion point and has valuation \(e\); the higher
terms of the formal group law have strictly larger valuation.
Inertia over \(L\) fixes this map, because it fixes \(\pi_L\)
and the algebraically closed residue field. Thus all \(p\)-torsion
is \(L\)-rational. The chosen generator over \(K\) acts on the
two-dimensional leading-coefficient space by multiplication by
\[
 \lambda=\zeta^{c+e}\in\mathbf F_q^\times .
 \tag{B-RD69.9}
\]
If its \(\mathbf F_p\)-representation is reducible, an invariant
line gives \(\lambda\in\mathbf F_p^\times\); its action on the
entire module is then scalar. Otherwise the residual module
is irreducible over \(\mathbf F_p\).

Suppose the first polygon is one-slope and the second two-slope.
The latter has a stable high line, so their common residual
module is reducible, hence scalar by the preceding paragraph.
In particular every line is stable under the actual inertia
action, and the second first-level module is \(L\)-rational.
Write its parameter as \(B\); then \(p\mid B\).
Transport its high line across the actual torsion isomorphism
to a line in the first group, and quotient both groups.
Lemma~\ref{lem:B-RD69-quotients} gives an actual \(p^2\)-comparison
with first parameter \(e\), prime to \(p\). If \(B<e\), the
second parameter is \(pB\), divisible by \(p\). Such a comparison
is impossible by Lemma~\ref{lem:B-RD69-rules}: the first parameter
forces nontrivial residual wild inertia in the two quotients.
If \(B>e\), the second parameter is \(s-B\), prime to \(p\)
because \(p\nmid s\). Equality of the two parameters forces
\(s-B=e\), or \(B=pe\), contrary to the strict two-slope bound.
The equality \(B=e\) cannot occur here, because \(p\mid B\)
whereas \(p\nmid e\). This excludes mixed polygons.

If both polygons have one slope and their common residual
module is irreducible, Theorem~\ref{thm:B-RD67} already recovers
their special characters from their actual \(p^2\)-torsion.
If it is reducible, both residual actions are scalar.
An isomorphism of the actual residual representations equates
the two scalars in (B-RD69.9). With the same generator and the
same \(e\), this gives \(c_1=c_2\pmod N\), hence the same
special character. This uses the actual scalar action, not
merely its determinant.
\end{proof}

\begin{proposition}[Two unequal polygons below \(p^2\)]
\label{prop:B-RD69-two-slope}
\status{Candidate; complete proof supplied, independent audits pending.}
Suppose in addition \(e<p^2\). Two whole potentially
supersingular height-two groups with isomorphic actual
\(p^3\)-torsion and two-slope first-level polygons have the
same finite special characters, regardless of the square
class of their common finite determinant.
\end{proposition}
\begin{proof}
Write their special tangent exponents as \(c_i\).
Their generic determinant is \(\chi_p\epsilon_i\), with
\(\epsilon_i\) finite and \(\mu_{p-1}\)-valued.
The actual torsion comparison and \(\mu_p\subset K\) imply
\(\epsilon_1=\epsilon_2\); reduction modulo \(p\) is already
injective on these finite characters. Its special reduced-norm
description gives
\[
 c_1\equiv c_2\pmod{p-1}.
 \tag{B-RD69.10}
\]
This equality, rather than a normalization to a primitive norm,
is all that is needed below.

If residual wild inertia is nontrivial, the high line is common.
Lemma~\ref{lem:B-RD69-rules} handles every case in which at
least one \(B_i\) is prime to \(p\). This expressly includes
the possible critical value \(B_i=e\).
For the remaining cases both \(B_i\) are divisible by \(p\),
so neither equals \(e\). If both are below \(e\), write
\(B_i=pb_i\). Their common high-line character and (B-RD69.10)
give \(b_1\equiv b_2\pmod{p-1}\). But
\[
 0<b_i<e/p<p,
 \tag{B-RD69.11}
\]
so both are in the interval \(\{1,\ldots,p-1\}\).
Reduction modulo \(p-1\) is injective on this interval, and
\(b_1=b_2\). The common high-line character identifies \(c_1=c_2\).
If both \(B_i>e\), their high quotients have parameters
\(s-B_i\) prime to \(p\); their actual \(p^2\)-comparison and
Lemma~\ref{lem:B-RD69-rules} give equal quotient characters,
hence equal original ones.
If the two \(B_i\)'s lie on opposite sides of \(e\), the high
quotient parameters have opposite \(p\)-divisibility and their
actual \(p^2\)-comparison is impossible by the same lemma.

It remains that the residual module is tame. Write \(B_i=pb_i\),
where \(0<b_i<e\); again \(B_i=e\) is impossible.
Two \(b_i\) prime to \(p\) are settled by the lemma.
Suppose \(p\nmid b_1=b\) and \(b_2=pk\).
Then \(B_2=p^2k>e\), since \(e<p^2\).
Choose a common residual vector outside both high lines and
corresponding primitive lifts. Their common point field has
degree \(p^2\) over \(L\), forced by the first valuation
\(b/p^2\); its two normalized coordinate valuations are
\(b\) and \(pk\). If the high lines differ, choose a translation
high in the first model and low in the second.
The prime-to-\(p\) valuation of the first coordinate fixes its
exact lower break as \(p^2\alpha_1-b\). The universal lower
bound for the second coordinate would give
\[
 p^2(\alpha_1-pk)\le b-pk .
 \tag{B-RD69.12}
\]
However \(\alpha_1=s-pb\ge e+p\), since \(b<e\) is integral,
while \(pk<e\) and \(b<e<p^2\). The left side is greater than
\(p^3\) and the right side less than \(p^2\). Thus the high
lines agree. Quotient by that common line.
The second new parameter \(s-B_2\) is prime to \(p\).
If \(B_1<e\), the first new parameter \(pB_1\) is divisible
by \(p\), an impossibility. If \(B_1>e\), equality of the two
prime-to-\(p\) new parameters forces \(B_1=B_2\), contrary
to \(p\nmid b\) and \(p\mid b_2\).

Finally suppose \(b_i=pk_i\) for both groups.
Then \(B_i=p^2k_i>e\). If their high lines agree, quotienting
by that line gives two prime-to-\(p\) parameters and recovers
the characters. If their high lines differ, quotient the
first by its high line and the second by the corresponding
stable low line. Both lines are stable because they correspond
under the actual residual isomorphism.
The new parameters \(s-B_1\) and \(B_2/p=pk_2\) have opposite
\(p\)-divisibility, contradicting their inherited actual
\(p^2\)-torsion isomorphism.
These cases exhaust the residual actions and all divisibilities.
Only \(p\)-divisible original parameters are quotiented in this
case analysis, so the critical quotient at \(B=e\) is never
silently treated as a two-slope quotient.
\end{proof}

\begin{theorem}[The first tame boundary has exact index two]
\label{thm:B-RD69}
\status{Candidate; complete proof supplied, independent audits pending.}
For every prime \(p\ge5\), in the full B.4 monodromy setting
with height two, perfect residue field and \(\mu_p\subset K\),
\[
 \boxed{\displaystyle \nu_{\min}(p,p^2-1,2)=2.}
 \tag{B-RD69.13}
\]
In particular \(\nu_{\min}(5,24,2)=2\).
More generally, actual \(p^3\)-torsion determines the requested
character whenever \(p-1\mid e\), \(p\nmid e\), and \(e<p^2\).
\end{theorem}
\begin{proof}
Pass to the completed maximal unramified base, retaining the
absolute ramification and the actual torsion comparison.
The determinant of an allowed height-two monodromy system is
\(\chi_p^d\epsilon\), with \(d\in\{0,1,2\}\) and finite
\(\mu_{p-1}\)-valued \(\epsilon\), by the dimension and
filtration argument in Theorem~\ref{thm:B-RD68}.
As \(p\nmid e\), the cyclotomic image modulo \(p^3\) has
order \(p^2\). Equality of the actual determinants gives the
same \(d\), since \(|d_1-d_2|\le2<p\), and then the same
\(\epsilon\).
Dimensions zero and two are determined by residual inertia
and Cartier duality. In total dimension one, two systems
filtered by potentially good height-one pieces have their
finite characters determined modulo \(p\).
A comparison of a whole potentially supersingular group
with a filtered one is excluded by
Proposition~\ref{prop:B-RD63-mixed}, because \(e<p^2<p^3\).
For the two whole groups, the possible Newton-polygon pairs
are exhausted by Propositions~\ref{prop:B-RD69-one-slope}
and~\ref{prop:B-RD69-two-slope}. They prove the asserted
upper bound throughout the stated range.

At \(e=p^2-1>p\), the explicit comparison of
Theorem~\ref{thm:B-RD63-arc}, with \(n=2\), gives two permitted
height-two monodromy systems with equal actual \(p^2\)-torsion
and different requested characters. Thus manuscript index one
does not suffice, while index two does. This proves the exact
value, not only a mixed-case necessary bound.
\end{proof}

\begin{remark}[Scope and the critical value]
\label{rem:B-RD69-scope}
\status{Candidate scope statement; independent audits pending.}
The one-slope/two-slope exclusion is valid for every
prime-to-\(p\) absolute index divisible by \(p-1\).
The recovery proof for two unequal polygons uses \(e<p^2\)
in the small-parameter interval and in the tame divisibility
case split. It imposes no square-determinant hypothesis and
uses actual rank-\(p\) quotients, not transferred Hasse thickness.
The value \(B=e\), absent from the earlier odd-\(B\) normalized
proof, is included explicitly: it is prime to \(p\), is handled
before the quotient case analysis, and its actual high quotient
is correctly recorded as one-slope.
This result does not cover \(e=p(p-1)\), nor any other absolute
index divisible by \(p\).
The retained exact supplement
\path{code/reduction_descent/verify_one_slope_boundary.py}
reconstructs the Newton polygons with rational arithmetic and
counts quotient roots from the actual product formula.
For \(72\) sectors at \(p=5,7,11,13,17,19\), it checks
\(130{,}398\) unequal polygons and high quotients,
\(216\) one-slope product quotients, all \(72\) critical
\(B=e\) quotients, \(7{,}842\) low quotients,
\(2{,}099{,}574\) unrestricted tangent weights, and
\(76{,}764\) tame mixed-divisibility inequalities.
It also checks the one-/two-slope arithmetic at additional
indices above \(p^2\), where only that particular exclusion is
claimed. The author run passed in \(7.471134923995123\) seconds
at 21:28:53 UTC. The exact program, full output and pre-certificate
proof are retained. These finite checks supplement the all-prime
proofs; they do not establish existence or nonexistence of
formal-group comparisons by enumeration alone.
\end{remark}
